\setcounter{qcount}{0}

\begin{sleekbox}[palCoral]{Problem \#1}
    \begin{question}\leavevmode\par\medskip
       Find the angle between the vectors $\anglevector{2, 4, 6}$ and $\anglevector{1, 5, 4}$. Show work by hand and your verification on Mathematica.
    \end{question}\leavevmode\par\medskip
    \begin{solution}\leavevmode\par\medskip
    
    \end{solution}
    
\end{sleekbox}

\begin{sleekbox}[palCoral]{Problem \#2}
    \begin{question}\leavevmode\par\medskip
       Find the work done by the force $\Vec{F} = \anglevector{3, 1, 7}$; the object displaced by the force
undergoes a displacement $\Vec{d} = \anglevector{5, 3, 2}$.
    \end{question}\leavevmode\par\medskip
    \begin{solution}\leavevmode\par\medskip
    
    \end{solution}
    
\end{sleekbox}

\begin{sleekbox}[palCoral]{Problem \#3}
    \begin{question}\leavevmode\par\medskip
       Find the force on an electron $\left(charge = q = 1.6 * 10^{-19} Coulombs\right)$ that is moving
    along the y axis with a speed of $8\times10^6 m/sec$ as it enters a magnetic field of strength 1.1
    tesla, directed in the positive z direction. (See example 6, page 841, section 13.4, for a
    similar problem in your book.)
    \end{question}\leavevmode\par\medskip
    \begin{solution}\leavevmode\par\medskip
    
    \end{solution}
    
\end{sleekbox}

\begin{sleekbox}[palCoral]{Problem \#4}
    \begin{question}\leavevmode\par\medskip
       Sketch the \textit{helix}: $\Vec{r}(t) = \anglevector{\cos(t), \sin(t),t}$
    \end{question}\leavevmode\par\medskip
    \begin{solution}\leavevmode\par\medskip
    
    \end{solution}
    
\end{sleekbox}

\begin{sleekbox}[palCoral]{Problem \#5}
    \begin{question}\leavevmode\par\medskip
       Write the parametric equations of a line that goes through the origin and is parallel to
    the vector $\anglevector{2, 4, 6}$. Then write the vector expression for the same line.
    \end{question}\leavevmode\par\medskip
    \begin{solution}\leavevmode\par\medskip
    
    \end{solution}
    
\end{sleekbox}

\begin{sleekbox}[palCoral]{Problem \#6}
    \begin{question}\leavevmode\par\medskip
       Write a vector expression for the curve which is given by $y=x \text{ and } z=x^2 + y^2$ using the parameter $t$. Then graph it on Mathematica.
    \end{question}\leavevmode\par\medskip
    \begin{solution}\leavevmode\par\medskip
    
    \end{solution}
    
\end{sleekbox}

\begin{sleekbox}[palCoral]{Problem \#7}
    \begin{question}\leavevmode\par\medskip
    For the function y = f(x), assuming f(x) is twice differentiable, this is another formula
    for curvature:
    \[
    k(x) = \frac{|f''(x)|}{\left( 1 + f'(x)^2 \right)^\frac{3}{2}}
    \]
        \begin{parts}
            \item Find the curvature for $f(x) = \ln(x), (x > 0)$ and find the point at which it is a maximum.
            \item What is the value of the maximum curvature?
            \item Find the curvature for $f(x) = e^x$ and find the point at which it is a maximum.
            \item What is the value of the maximum curvature?
            \item Compare your answers to parts b and d. What are the results? Why might this be happening?
        \end{parts}
    \end{question}\leavevmode\par\medskip
    \begin{solution}\leavevmode\par\medskip
        \begin{parts}
            \item 
            \item 
            \item 
            \item 
            \item 
            \item 
        \end{parts}
    \end{solution}
    
\end{sleekbox}

\newpage
